\section{Introduction: the eight questions behind a yes vote}
\label{sec:intro}

A bill does not pass because it is correct. It passes because enough members of
each chamber can answer, to their own satisfaction and to a constituent's, the
small set of questions that stand between a draft and a statute. This paper sets out
those questions for H. R. 9510 and answers each with evidence.

[DRAFTING INSTRUCTIONS: In the full stage, define for a lay legislator, in plain
language and before first substantive use, the following terms, drawing definitions
from the bill text in
\texttt{single-prompt-bill/auto-bill-02/final-bill/publication} (sections
\texttt{s2-findings} and \texttt{s3-amendment}) and the platform scope in
\texttt{physical-ai-oncology-trials/national-platform/new\_paper/final\_paper}
(\texttt{executive\_summary}): Physical AI; robot-patient interaction code; VVUQ;
verification before generation; the ten gates and the three hard catastrophe
predicates; ACCEPT, ESCALATE, and BLOCK; section 515D (21 U.S.C. 360e-5);
authorization versus appropriation; PAYGO and the House CutGo rule; and regular
order. State the three recurring passage questions and carry them as motifs:
``Would I defend this vote to a constituent who lost someone to a preventable
error?''; ``Can I explain the score to the Budget Committee?''; and ``Is there
daylight between my party and the other on this?'' Cite \cite{Kawchak2026HR9510Bill}
and \cite{Kawchak2026NationalPlatformPhysicalAI}.]

The map of the argument is one figure: the eight questions are an ordered path from
the mandate for a statute to the passage path itself. Figure 1 in the appendix
(\texttt{enactment/mermaid/diagrams/02-eight-passage-questions.md}) traces that path.

\begin{psgfig}
\node[psgone] (a) at (0,0) {Member receives\\H. R. 9510};
\node[psgdec] (q) at (4.0,0) {Eight\\questions};
\node[psgevid] (y) at (8.0,0) {A yes vote in\\both chambers};
\draw[psgedge] (a)--(q); \draw[psgedge] (q)--(y);
\end{psgfig}
\psgcaption{Figure I. The framework in miniature: the eight passage questions stand
between the bill on the desk and a yes vote in both chambers. [FULL STAGE: replace
with the colored TikZ reproduction of diagram 02.]}

This introduction will close, in the full stage, by naming the single objective:
not to win an argument, but to assemble a majority in the House and the votes to pass
the Senate.
